\section{方法与实验设置}

\subsection{整体流程}

本文采用统一流程完成模型比较：
原始数据先经过缺失值处理、异常值裁剪、标准化和分层划分，
随后在相同训练集、验证集和测试集上训练并评估不同模型。
验证集用于模型选择和超参数调优，
测试集仅用于最终性能报告。

\placeholder{补充整体 Pipeline 图。
建议展示 Dataset $\rightarrow$ Preprocessing $\rightarrow$
Train/Validation/Test Split $\rightarrow$ Model Training/Tuning
$\rightarrow$ Evaluation $\rightarrow$ Demo。}

\subsection{模型选择}

当前已完成 Logistic Regression 基准模型：
使用 \texttt{class\_weight="balanced"}、最大迭代次数 1000、
随机种子 42，并将其作为可解释的性能下限。
后续模型比较应在同一数据划分和同一预处理输出上完成。
表 \ref{tab:model_zoo} 给出了建议保留的模型类型及其作用。

\begin{table}[t]
  \caption{候选模型与选择动机}
  \label{tab:model_zoo}
  \centering
  \scriptsize
  \begin{tabular}{p{0.24\linewidth}p{0.37\linewidth}p{0.29\linewidth}}
    \toprule
    模型 & 选择动机 & 关键参数 \\
    \midrule
    Logistic Regression & 可解释基准，建立性能下限 & class\_weight, C \\
    KNN & 距离型非参数方法 & n\_neighbors, weights \\
    Decision Tree & 规则清晰，便于解释 & max\_depth, min\_samples\_split \\
    Random Forest & Bagging 集成，降低方差 & n\_estimators, max\_depth \\
    SVM & 间隔最大化分类器 & C, kernel, gamma \\
    XGBoost/MLP & 可选增强模型 & 按实际实现补充 \\
    \bottomrule
  \end{tabular}
\end{table}

\placeholder{根据实际完成情况补充各模型的训练代码和参数。
如果某个模型没有训练，不要在最终结果表中填写虚构指标。}

\subsection{超参数调优协议}

多模型阶段应使用验证集或交叉验证进行超参数选择。
若采用 GridSearchCV 或 RandomizedSearchCV，
需要报告搜索范围、交叉验证折数、选择指标和最优参数。
对于糖尿病筛查任务，调参时不宜只优化 Accuracy，
应同时关注 Recall、F1-score 和 ROC AUC。

\placeholder{补充实际调参方法和搜索空间。
例如 KNN 的邻居数、SVM 的 $C$ 和 \texttt{gamma}、
Random Forest 的树数量和深度等。}

\subsection{评价指标}

糖尿病预测是类别不平衡二分类问题，
单一 Accuracy 可能掩盖阳性样本漏判。
因此本文报告 Accuracy、Precision、Recall、F1-score 和 ROC AUC。
其中 Recall 衡量真实阳性样本被正确识别的比例；
在筛查场景中，假阴性意味着潜在高风险个体被漏掉，
因此 Recall 和 F1-score 比单纯 Accuracy 更需要关注。
ROC AUC 不依赖单一分类阈值，
用于衡量模型对阳性样本的概率排序能力。
